\documentclass[11pt]{article}
\usepackage[a4paper,margin=27mm]{geometry}
\usepackage{amsmath,amssymb,amsthm,booktabs,array}
\usepackage[colorlinks=true,linkcolor=blue,urlcolor=blue,citecolor=blue]{hyperref}
\newtheorem{theorem}{Theorem}
\newtheorem{corollary}[theorem]{Corollary}
\newcommand{\Q}{\mathbb{Q}}
\newcommand{\Res}{\operatorname{Res}}
\title{OEIS A386829, FORMULA 5, revisions 11--12\\
(2 November 2026): Mathar's recurrence\\
and an algebraic family}
\author{Haobo Ma\\\small ChronoAI Pte Ltd, Singapore
\and Wenlin Zhang\\\small National University of Singapore\\\small ChronoAI Pte Ltd, Singapore}
\date{7 November 2026}
\begin{document}
\maketitle
\begin{abstract}
We prove R. J. Mathar's recurrence and D-finiteness conjecture in
OEIS A386829, FORMULA 5, revision 11 (2 November 2026), approved as revision 12.
For Manyama's inclusive binomial sum, the ordinary generating function is
$\bigl(7/\sqrt{1-20z}-3\bigr)/(4+18z)$; its first-order
inhomogeneous differential equation gives the stated recurrence at every
$n\geq2$. A formal change of variable also proves the entry's alternative
formulas and gives a two-weight algebraic family. It also proves the
displayed recurrence conjectures in A386830 and A386831, and the
exponential, reversion and hypergeometric identities in related entries.
\end{abstract}

\section{The conjecture and its generating function}
Seiichi Manyama defined A386829 in 2025 by the finite sum
\begin{equation}\label{eq:definition}
 a_n=\sum_{k=0}^{n}3^k2^{n-k}
       \binom{2n+1}{k}\binom{2n-k}{n-k}\qquad(n\geq0).
\end{equation}
The summation includes both endpoints; in particular $a_0=1$ and $a_1=13$.
Mathar's conjecture in FORMULA 5 of the entry~\cite{oeis829}
is the recurrence in the following theorem, together with D-finiteness.
The history records its addition at 08:27:56 EST on 2 November 2026
(revision 11), followed by approval at 08:28:00 EST (revision 12).
All generating functions below are formal power series over $\Q$.
The square root denotes the unique series with constant coefficient one.

\begin{theorem}\label{thm:main}
For the sequence~\eqref{eq:definition}, put $A(z)=\sum_{n\geq0}a_nz^n$.
Then
\begin{equation}\label{eq:gf}
 A(z)=\frac{7(1-20z)^{-1/2}-3}{4+18z},
\end{equation}
and, for every integer $n\geq2$,
\begin{equation}\label{eq:recurrence}
 2n a_n+(20-31n)a_{n-1}+90(1-2n)a_{n-2}=0.
\end{equation}
Moreover, $A$ satisfies the nonzero homogeneous differential equation
\begin{equation}\label{eq:ode}
 (2-31z-180z^2)A''-(42+630z)A'-270A=0.
\end{equation}
\end{theorem}

We first establish a coefficient identity which also explains the
alternative formulas in the OEIS entry. The all-$r$, two-weight
generalization is not posed in the cited entries. For an integer
$r\geq1$ and $u,v\in\Q$ define
\begin{equation}\label{eq:family}
 d^{(r)}_n(u,v)=\sum_{k=0}^{n}u^kv^{n-k}
 \binom{(r+1)n+1}{k}\binom{(r+1)n-k}{n-k}.
\end{equation}
Write $p=u+v$ and $q=u$. Expanding the negative binomial series gives
\begin{align}
 d^{(r)}_n(u,v)
 &= [x^n]\frac{(1+ux)^{(r+1)n+1}}{(1-vx)^{rn+1}}\label{eq:coef1}\\
 &= [y^n]\frac{1}{(1-qy)(1-py)^{rn+1}}\label{eq:coef2}\\
 &= \sum_{k=0}^{n}p^kq^{n-k}\binom{rn+k}{k}.
 \label{eq:sum2}
\end{align}
Here is a direct justification of the nontrivial change from
\eqref{eq:coef1} to~\eqref{eq:coef2}. Write $[x^n]F(x)$ as the
formal residue $\Res_{x=0} F(x)x^{-n-1}\,dx$, and substitute
$x=y/(1-uy)$. The four contributions from the numerator, denominator,
$x^{-n-1}$ and $dx$ have respective powers
$-(r+1)n-1$, $rn+1$, $n+1$ and $-2$ of $1-uy$.
Their sum is $-1$, giving exactly~\eqref{eq:coef2}.
Formal residues are invariant under a substitution with leading term $y$:
this follows on Laurent monomials from the vanishing residue of a
derivative, and from $\Res x'(y)/x(y)=1$.
Every coefficient calculation here is finite.

The further substitution $y=t/(1+pt)$ in~\eqref{eq:coef2} gives
\begin{equation}\label{eq:sum3}
 d^{(r)}_n(u,v)
 =[t^n]\frac{(1+pt)^{(r+1)n+1}}{1+vt}
 =\sum_{k=0}^{n}p^k(-v)^{n-k}\binom{(r+1)n+1}{k}.
\end{equation}
Indeed, the resulting power of $1+pt$ in the residue is $(r+1)n+1$.
Thus these equalities prove all four displayed alternative formulas
for A386829, rather than assuming any one of them as its definition.

\section{A formal coefficient formula and the proof}
We use a standard form of Lagrange inversion, also recorded as
Theorem 2.1.1, equation (2.1.5), in Gessel~\cite{gessel}.
If $\Phi(0)=1$, $T=z\Phi(T)$, and $H$ is a power series, then
\begin{equation}\label{eq:lagrange}
 \sum_{n\geq0}z^n[t^n]H(t)\Phi(t)^n
 =\frac{H(T)}{1-T\Phi'(T)/\Phi(T)}.
\end{equation}
For completeness, substitute $z=t/\Phi(t)$ in the residue for the
$n$th coefficient on the right. Its derivative contributes
$(1-t\Phi'(t)/\Phi(t))/\Phi(t)$; cancellation leaves
$\Res H(t)\Phi(t)^nt^{-n-1}\,dt$, the coefficient on the left.
The solution $T\in z\Q[[z]]$ exists uniquely by successive determination
of its coefficients, so this substitution is valid also at $n=0$.

Taking $\Phi(t)=(1-pt)^{-r}$ and
$H(t)=1/((1-qt)(1-pt))$ in~\eqref{eq:lagrange} proves
\begin{equation}\label{eq:general-gf}
 D_r(z):=\sum_{n\geq0}d^{(r)}_n(u,v)z^n
 =\frac{1}{(1-qT)(1-(r+1)pT)},\qquad
 T(1-pT)^r=z.
\end{equation}
Every denominator on the right has constant coefficient one.
For $r=1$, $p=5$ and $q=3$, put $s=\sqrt{1-20z}$.
The solution with zero constant coefficient is $T=(1-s)/10$, hence
\[
 D_1(z)=\frac{10}{s(7+3s)}
        =\frac{7/s-3}{4+18z},
\]
which proves~\eqref{eq:gf} for the literal sum~\eqref{eq:definition}.

\begin{proof}[Completion of the proof of Theorem~\ref{thm:main}]
Let $C=(1-20z)^{-1/2}$. Then $(1-20z)C'=10C$ and
$(4+18z)A=7C-3$. Differentiating the latter equality and eliminating
$C$ gives
\begin{equation}\label{eq:inhom}
 (2-31z-180z^2)A'-(11+270z)A=15.
\end{equation}
For $n\geq2$, its coefficient of $z^{n-1}$ is
\[
 2na_n-\bigl(31(n-1)+11\bigr)a_{n-1}
        -\bigl(180(n-2)+270\bigr)a_{n-2}=0,
\]
which is~\eqref{eq:recurrence}. The constant coefficient instead gives
$2a_1-11a_0=15$, consistent with $a_0=1,a_1=13$.
Differentiating~\eqref{eq:inhom} gives~\eqref{eq:ode}.
Its leading polynomial has constant coefficient $2$, so it is a nonzero
homogeneous polynomial differential equation. This proves D-finiteness
with its usual meaning~\cite[\S1]{stanley}.
\end{proof}

In particular the square-root representation yields the algebraic identity
\begin{equation}\label{eq:algebraic}
 (1-20z)\bigl((4+18z)A(z)+3\bigr)^2=49.
\end{equation}
No assertion about analytic convergence is needed for these arguments.

\section{Directly related sequences}
The same calculation gives the following consequences for the four
entries directly cross-referenced by A386829.

\begin{corollary}\label{cor:siblings}
Let $C=(1-20z)^{-1/2}$. The generating functions of
A385669 and A386763 are respectively
\begin{equation}\label{eq:siblings}
 B(z)=\frac{4C-1}{3+4z},\qquad
 E(z)=\frac{(1+15z)C+1}{2+9z}.
\end{equation}
The coefficients $b_n$ of $B$ satisfy $b_0=1,b_1=12$ and
\begin{equation}\label{eq:siblingrec}
 3n b_n+(30-56n)b_{n-1}-40(2n-1)b_{n-2}=0
 \qquad(n\geq2).
\end{equation}
The generating functions for A386830 and A386831 are the algebraic
series $D_2$ and $D_3$ in~\eqref{eq:general-gf}, with $p=5,q=3$.
Both are D-finite.
\end{corollary}
\begin{proof}
For A385669~\cite{oeis669}, use $(u,v)=(2,3)$ and $r=1$ in
\eqref{eq:family}--\eqref{eq:general-gf}. Rationalizing with
$T=(1-s)/10$ gives $B$ in~\eqref{eq:siblings}.
The identities $(3+4z)B=4C-1$ and $(1-20z)C'=10C$ give
\[
 (3-56z-80z^2)B'-(26+120z)B=10.
\]
Taking the coefficient of $z^{n-1}$ gives~\eqref{eq:siblingrec}.

The definition of A386763~\cite{oeis763} is
$e_n=\sum_{k=0}^n5^k3^{n-k}\binom{n+k-1}{k}$, with $e_0=1$.
Here the coefficient at $n=k=0$ is understood as
$[x^0](1-5x)^0=1$.
Thus $e_n=[x^n](1-3x)^{-1}(1-5x)^{-n}$ for all $n\geq0$.
Use $\Phi=(1-5t)^{-1}$ and $H=(1-3t)^{-1}$ in
\eqref{eq:lagrange}. The result is $E=(1-5T)A$.
Substitution of $T=(1-s)/10$ and $s^2=1-20z$ yields the displayed
formula for $E$.

The defining sums of A386830 and A386831~\cite{oeis830,oeis831}
are exactly~\eqref{eq:family} with $(u,v)=(3,2)$ and $r=2,3$.
Equation~\eqref{eq:general-gf} proves their algebraicity.
To see D-finiteness directly, the field $\Q(z,T)$ is a finite-dimensional
extension of $\Q(z)$, of degree at most $r+1$. Implicit differentiation
of $T(1-5T)^r=z$ gives
\[
 T'=\frac{1}{(1-5T)^{r-1}(1-5(r+1)T)}\in\Q(z,T).
\]
Consequently differentiation preserves this field. The derivatives
of $D_r$ are linearly dependent over $\Q(z)$; clearing denominators
gives a nonzero homogeneous polynomial differential equation.
This is the standard algebraic-to-D-finite implication
\cite[Theorem 2.1]{stanley}, with its hypotheses explicit here.
\end{proof}

\section{Two further recurrence conjectures}
Let $d_n=d^{(r)}_n(3,2)$ and $D=D_r$. The entries A386830 and A386831
pose specific recurrences, not just D-finiteness. Both follow from
\eqref{eq:general-gf}.

\begin{corollary}\label{cor:recurrences}
For A386830, Mathar's conjecture in FORMULA 5, revision 9
(19 August 2025; approved as revision 10), holds for every $n\geq2$:
\begin{multline}\label{eq:r2}
8n(2n-3)d_n+6(-108n^2+207n-80)d_{n-1}\\
 +405(3n-2)(3n-4)d_{n-2}=0.
\end{multline}
For A386831, Mathar's conjecture in FORMULA 5, revision 8
(2 November 2026; approved as revision 9), holds for every $n\geq4$:
\begin{equation}\label{eq:r3}
 \sum_{j=0}^4 P_j(n)d_{n-j}=0,
\end{equation}
where
\begin{align*}
P_0(n)&=18744n(3n-1)(3n-2),\\
P_1(n)&=203878391n^3-1040550807n^2+1658851582n-822702720,\\
P_2(n)&=27(-252935041n^3+1277608165n^2\\
 &\hspace{36mm}-1931547954n+777640080),\\
P_3(n)&=3645(-38797609n^3+259381745n^2\\
 &\hspace{36mm}-614869982n+523823440),\\
P_4(n)&=-16309333800(2n-7)(4n-13)(4n-11).
\end{align*}
\end{corollary}
\begin{proof}
Put $Z=t(1-5t)^r$, $V=1-3t$, $W=1-5(r+1)t$, and $K=1-5t$.
Under $z=Z$, the Euler derivative $\theta=z\,d/dz$ becomes
\[
 \mathcal L=\frac{tK}{W}\frac{d}{dt},\qquad D(Z)=\frac1{VW}.
\]
Here is an explicit polynomial identity proving each recurrence.
Define $N_0=1$ and, for $k\geq0$,
\begin{equation}\label{eq:certificate}
N_{k+1}=tK\{VW N'_k+((k+1)3W+(2k+1)5(r+1)V)N_k\}.
\end{equation}
Differentiation proves
$\mathcal L^k(1/(VW))=N_k/(V^{k+1}W^{2k+1})$.
Write $P_j(X+j)=\sum_{k=0}^m c_{jk}X^k$. For $r=2$, use the
three coefficient polynomials in~\eqref{eq:r2} and $m=2$; for $r=3$,
use those in~\eqref{eq:r3} and $m=3$. Expansion gives the following
identities in $\mathbb Z[t]$:
\begin{equation}\label{eq:zero-polynomial}
\sum_{j,k}c_{jk}Z^jN_kV^{m-k}W^{2(m-k)}
 =R_r(Z)V^{m+1}W^{2m+1},
\end{equation}
where
\[
 R_2(z)=-30z,\qquad
 R_3(z)=338670z-181329840z^2+604049400z^3.
\]
This equality is a finite polynomial identity: the displayed recurrence
for $N_k$, the listed $P_j$, and ordinary multiplication specify every
coefficient, with no unknown series. The supplied
\texttt{verify\_polynomials.py} expands both sides using integer arithmetic.
Since $V(0)=W(0)=1$, division in~\eqref{eq:zero-polynomial} is valid and gives
\[
 \sum_j z^jP_j(\theta+j)D(z)=R_r(z).
\]
The coefficient of $z^n$ on the left is exactly
$\sum_jP_j(n)d_{n-j}$, with negative subscripts omitted. Since
$\deg R_2<2$ and $\deg R_3<4$, this proves the stated onsets.
In particular, the lower-index terms have not been silently discarded.
\end{proof}

\section{Shifted sums, reversions and secondary entries}
The same argument covers more than the four immediate cross-references.
For $r\geq1$, $u,v\in\Q$, $p=u+v$, $q=u$, define $e_0=1$ and
\[
 e_n=\sum_{k=0}^n p^kq^{n-k}\binom{rn+k-1}{k}\quad(n\geq1),
 \qquad \Phi(x)=\frac{(1+ux)^{r+1}}{(1-vx)^r}.
\]
\begin{corollary}\label{cor:shifted}
With $T(1-pT)^r=z$ and $T(0)=0$, the following identities hold:
\begin{align}
 e_n&=[x^n]\Phi(x)^n
 =\sum_{k=0}^np^k(-v)^{n-k}\binom{(r+1)n}{k},\label{eq:eshift}\\
 E_r(z)&=\sum_{n\geq0}e_nz^n
 =\frac{1-pT}{(1-qT)(1-(r+1)pT)},\label{eq:egeneral}\\
 \exp\left(\sum_{n\geq1}\frac{e_n}{n}z^n\right)
 &=\frac{R(z)}z=\frac{1}{(1-qT)(1-pT)^r},
 \quad R=z\Phi(R).\label{eq:exp}
\end{align}
The coefficients of $R/z$ are integers if $u,v$ are integers.
For $n\geq1$ and $q\ne0$, $p/q\ne1$, the hypergeometric identity is
\begin{multline}\label{eq:hypergeom}
 e_n=q^n(1-p/q)^{-rn}\\
 -p^n\binom{(r+1)n-1}{n}
 \left({}_2F_1(1,(r+1)n;n+1;p/q)-1\right).
\end{multline}
The hypergeometric function here is a rational function of its argument,
continued from its series at zero.
\end{corollary}
\begin{proof}
The substitutions used in~\eqref{eq:coef1}--\eqref{eq:sum3}, with the
numerator exponent reduced by one, give
$e_n=[y^n](1-qy)^{-1}(1-py)^{-rn}$ and~\eqref{eq:eshift}.
They are valid at $n=0$ with the displayed convention $e_0=1$.
Formula~\eqref{eq:lagrange} now gives~\eqref{eq:egeneral}.
Substitution shows $R=T/(1-qT)$ and $R=z\Phi(R)$, proving the
second equality in~\eqref{eq:exp}. For the first, Lagrange inversion gives
\[
 [z^n]\log\Phi(R)
 =\frac1n[x^{n-1}]\frac{\Phi'(x)}{\Phi(x)}\Phi(x)^n
 =\frac1n[x^n]\Phi(x)^n=\frac{e_n}{n}.
\]
All logarithms have constant coefficient zero. Exponentiation proves
\eqref{eq:exp}. Integer coefficients follow recursively from
$R=z\Phi(R)$, since $\Phi\in\mathbb Z[[x]]$ for integral $u,v$.

For~\eqref{eq:hypergeom}, put $M=(r+1)n$. The identity
\[
 \binom{M-1}{n}\frac{(M)_k}{(n+1)_k}
 =\binom{M+k-1}{n+k}
\]
identifies the hypergeometric tail with the tail of
$(1-x)^{-rn}=\sum_{j\geq0}\binom{rn+j-1}{j}x^j$ after $j=n$.
Subtract the finite initial segment and multiply by $q^n$, with $x=p/q$.
This proves the formula first near zero. Its rational continuation is
explicit from Euler's transformation~\cite[15.8.1]{dlmf}:
\[
 {}_2F_1(1,M;n+1;x)
 =(1-x)^{-rn}\sum_{k=0}^{rn-1}
       \frac{n}{n+k}(-1)^k\binom{rn-1}{k}x^k.
\]
Thus no branch choice beyond rational continuation is involved, and
evaluation at $x=5/3$ is legitimate.
\end{proof}

For $(u,v)=(3,2)$ and $r=1,2,3$, these are precisely the coefficient,
exponential/reversion and Spezia hypergeometric formulas in A386763,
A386764 and A386765~\cite{oeis763,secondary}. For $(u,v)=(2,3)$,
the family $D_r$ includes A385669, A385670 and A385671; $E_1$ is A386766.
Their source identities follow from the same substitutions.

Two further displayed source recurrences follow by direct differentiation:
if $f_n$ is A386766, then, for $n\geq3$,
\begin{multline}\label{eq:766}
9nf_n+6(-18n-5)f_{n-1}+80(-17n+37)f_{n-2}\\
 +800(-2n+5)f_{n-3}=0.
\end{multline}
Indeed, use~\eqref{eq:certificate} with $r=1$, $q=2$, $V=1-2t$,
replace $3W$ by $2W$, and set $N_0=1-5t$. For the coefficient
polynomials in~\eqref{eq:766}, identity~\eqref{eq:zero-polynomial}
has $m=1$ and $R(z)=-75z+300z^2$.
If $w_n=[z^n]R(z)/z$ for $(u,v)=(3,2),r=1$, then $w_n$ is A386769 and
\begin{equation}\label{eq:769}
 2(n+1)w_n+(-31n+29)w_{n-1}+90(-2n+1)w_{n-2}=0\quad(n\geq2).
\end{equation}
To verify this, substitute $F=1/((1-3t)(1-5t))$ and $z=t(1-5t)$ in
\[
 (2-31z-180z^2)\theta F+(2-2z-270z^2)F=2+30z;
\]
clearing denominators gives a polynomial identity. Coefficient
extraction proves~\eqref{eq:769}. Finally Lagrange inversion gives
$w_n=(n+1)^{-1}[x^n]\Phi(x)^{n+1}$, which proves both of the
coefficient formulas in A386769.

For the secondary entry A383888 take $(u,v)=(1,2),r=1$ in
Corollary~\ref{cor:shifted}. It gives every displayed finite-sum,
coefficient, exponential and hypergeometric identity there. It also
gives Kotesovec's stated asymptotic. In fact, with $s=\sqrt{1-12z}$,
\[
 E_1(z)=\frac{3(1+s)}{s(5+s)}
       =\frac{3}{5s}+\frac{12}{5(5+s)}.
\]
The square root with $s(0)=1$ is analytic on the plane slit along
$[1/12,\infty)$ and has positive real part there, so $5+s$ does not vanish.
At the unique dominant singularity $z=1/12$, the second summand has
a convergent expansion in $s$, whose first nonanalytic term is a
constant times $s$. The binomial coefficient asymptotic (or the
transfer theorem~\cite[Theorems VI.1 and VI.3]{flajolet}) therefore gives
\[
 [z^n]E_1(z)\sim\frac{3\,12^n}{5\sqrt{\pi n}},
\]
exactly the FORMULA contribution dated 4 August 2025.

\section{Coverage and verification scope}
\begin{center}
\small
\begin{tabular}{@{}p{0.30\linewidth}p{0.61\linewidth}@{}}
\toprule
Source item & Mathematical coverage\\
\midrule
A386829, NAME & Proved: literal inclusive sum; Theorem~\ref{thm:main}.\\
A386829, FORMULA 1--4 & Proved: all four identities in
\eqref{eq:coef1}--\eqref{eq:sum3}, with $r=1,u=3,v=2$.\\
A386829, Mathar's conjecture, FORMULA 5 & Proved: recurrence for every $n\geq2$,
and D-finiteness; Theorem~\ref{thm:main}.\\
A385669 & Proved: four source identities. Explicit generating function
and recurrence: not posed; proved in Corollary~\ref{cor:siblings}.\\
A386830, FORMULA 1--5 & Proved: four identities, D-finiteness and the
conjectured recurrence for $n\geq2$; Corollary~\ref{cor:recurrences}.\\
A386831, FORMULA 1--5 & Proved: four identities, D-finiteness and the
conjectured recurrence for $n\geq4$; Corollary~\ref{cor:recurrences}.\\
A386763; A386764--765 & Proved: all five source formulas, including
Spezia's hypergeometric representations ($n\geq1$);
Corollary~\ref{cor:shifted}.\\
A385670--671 & Proved: the four source identities. Algebraicity for
the swapped weights: not posed; proved in \eqref{eq:general-gf}.\\
A386766 & Proved: all source formulas, including Mathar's stated
recurrence ($n\geq3$); \eqref{eq:766}.\\
A386769 & Proved: the two coefficient formulas and Mathar's stated
recurrence ($n\geq2$); \eqref{eq:exp}, \eqref{eq:769}.\\
A383888 & Proved: its five identities and Kotesovec's asymptotic;
Corollary~\ref{cor:shifted} and the displayed singular expansion.\\
General parameter family and explicit generating functions & Not posed;
proved in \eqref{eq:general-gf} and Corollary~\ref{cor:shifted}.\\
Other entries reached only by further cross-references & Not covered;
the table specifies the twelve entries treated here.\\
\bottomrule
\end{tabular}
\end{center}

The declaration
\href{https://github.com/the-omega-institute/trureturing/blob/074df09ddbbef1843ddb383577bfaa827f753e36/D5/S3/ArithSums/MatharA386829.lean#L233}{\texttt{D5.S3.ArithSums.MatharA386829.result}}
states the full recurrence~\eqref{eq:recurrence} for the literal sum,
together with existence of a nonzero homogeneous differential equation
for its rational formal generating series. The same file's
\texttt{homogeneous\_equation} gives~\eqref{eq:ode}.
The closed generating function, its parameter extension and all sibling
statements have written proofs above. In particular, the pinned declaration
does not cover \eqref{eq:r2} or \eqref{eq:r3}.
The supplied \texttt{verify.py}
checks exact integer coefficients through $n=80$ for the main sum,
the recurrence, the differential equations and~\eqref{eq:algebraic};
it also checks nine choices of $(r,u,v)$ through $n=30$, their
shifted and reversion identities, the hypergeometric formulas for
$1\leq n\leq30$, and the added recurrences through $n=30$.
These finite coefficient checks are distinct from the proofs for all indices.
The separate \texttt{verify\_polynomials.py} verifies the full polynomial
identities in the recurrence proofs, not a sample of their values.

\begin{thebibliography}{99}
\bibitem{oeis829} S. Manyama, OEIS A386829, NAME and FORMULA,
revisions 11--12, 2 November 2026; recurrence conjecture contributed by
R. J. Mathar in revision 11 on 2 November 2026.
\url{https://oeis.org/A386829}.
\bibitem{oeis669} S. Manyama, OEIS A385669, NAME and FORMULA,
revision 13, 4 August 2025. \url{https://oeis.org/A385669}.
\bibitem{oeis763} S. Manyama, OEIS A386763, NAME and FORMULA,
revision 19, 3 August 2025. \url{https://oeis.org/A386763}.
\bibitem{oeis830} S. Manyama, OEIS A386830, NAME and FORMULA,
revision 10, 19 August 2025; recurrence conjecture by R. J. Mathar.
\url{https://oeis.org/A386830}.
\bibitem{oeis831} S. Manyama, OEIS A386831, NAME and FORMULA,
revision 9, 2 November 2026; recurrence conjecture by R. J. Mathar.
\url{https://oeis.org/A386831}.
\bibitem{secondary} S. Manyama, secondary OEIS entries, NAME and FORMULA:
\href{https://oeis.org/A385670}{A385670}, revision 11, 4 August 2025;
\href{https://oeis.org/A385671}{A385671}, revision 13, 4 August 2025;
\href{https://oeis.org/A386764}{A386764}, revision 21, 3 August 2025;
\href{https://oeis.org/A386765}{A386765}, revision 21, 3 August 2025;
\href{https://oeis.org/A386766}{A386766}, revision 15, 3 August 2025;
\href{https://oeis.org/A386769}{A386769}, revision 14, 3 August 2025;
\href{https://oeis.org/A383888}{A383888}, revision 38, 4 August 2025.
The hypergeometric representations are due to S. Spezia (2--3 August 2025),
the recurrences to R. J. Mathar (3 August 2025), and the A383888 asymptotic
to V. Kotesovec (4 August 2025).
\bibitem{gessel} I. M. Gessel, Lagrange inversion,
\emph{J. Combin. Theory Ser. A} \textbf{144} (2016), 212--249.
\href{https://doi.org/10.1016/j.jcta.2016.06.018}{doi:10.1016/j.jcta.2016.06.018};
\href{https://arxiv.org/abs/1609.05988}{arXiv:1609.05988v1},
Theorem 2.1.1, equation (2.1.5).
\bibitem{stanley} R. P. Stanley, Differentiably finite power series,
\emph{European J. Combin.} \textbf{1} (1980), 175--188,
Theorem 2.1, p.~178.
\href{https://doi.org/10.1016/S0195-6698(80)80051-5}{doi:10.1016/S0195-6698(80)80051-5}.
\bibitem{flajolet} P. Flajolet and R. Sedgewick,
\emph{Analytic Combinatorics}, Cambridge University Press, 2009,
Theorems VI.1 and VI.3.
\href{https://ac.cs.princeton.edu/home/}{Author book site}.
\bibitem{dlmf} NIST Digital Library of Mathematical Functions,
\S15.8(i), equation (15.8.1), Euler transformation.
\href{https://dlmf.nist.gov/15.8.E1}{https://dlmf.nist.gov/15.8.E1}.
\end{thebibliography}
\bigskip\noindent\small
\copyright\ 2026 Haobo Ma and Wenlin Zhang.
Licensed under Creative Commons Attribution 4.0 International (CC BY 4.0).
\end{document}
